\documentclass[11pt,a4paper]{article}

\usepackage[margin=2.5cm]{geometry}
\usepackage{amsmath,amssymb,amsthm}
\usepackage{bm}
\usepackage{booktabs}
\usepackage{xcolor}
\usepackage{hyperref}
\usepackage{enumitem}

\newtheorem{theorem}{Theorem}[section]
\newtheorem{lemma}[theorem]{Lemma}
\newtheorem{corollary}[theorem]{Corollary}
\newtheorem{proposition}[theorem]{Proposition}
\newtheorem{remark}[theorem]{Remark}
\newtheorem{assumption}[theorem]{Assumption}

\theoremstyle{definition}
\newtheorem{definition}[theorem]{Definition}

\newcommand{\R}{\mathbb{R}}
\newcommand{\dom}{\Omega}
\newcommand{\pore}[1]{\Omega_{#1}}
\newcommand{\iface}{\gamma}
\newcommand{\vel}{\bm{u}}
\newcommand{\pres}{p}
\newcommand{\stress}{\bm{\sigma}}
\newcommand{\normal}{\bm{n}}
\newcommand{\trac}{\bm{\lambda}}
\newcommand{\DtN}{\mathbf{G}}
\newcommand{\Schur}{\mathbf{S}}
\newcommand{\Shat}{\widehat{\Schur}}
\newcommand{\Shatuu}{\widehat{\Schur}{}^{uu}}
\newcommand{\Fhat}{\widehat{\bm{F}}}
\newcommand{\one}{\mathbf{1}}
\newcommand{\zero}{\bm{0}}
\newcommand{\Range}{\operatorname{Range}}
\newcommand{\Ker}{\operatorname{Ker}}
\newcommand{\rank}{\operatorname{rank}}
\newcommand{\modeset}{\mathcal{M}}
\newcommand{\loadvec}{\bm{f}}
\newcommand{\Vspace}{\mathcal{V}}
\newcommand{\Wspace}{\mathcal{W}}
\newcommand{\Lspace}{\Lambda}
\newcommand{\Pcal}{\mathcal{P}}

\hypersetup{colorlinks=true,linkcolor=blue,urlcolor=blue,citecolor=blue}

\begin{document}

\title{\textbf{Theory of the Affine-DDPNM}\\[4pt]
       \large Best Approximation, Nested Decomposition, and Error Separation\\[4pt]
       \large Version 2 --- revised 2026-08-07}
\author{}
\date{}
\maketitle

\begin{center}
\fbox{\begin{minipage}{0.92\textwidth}
This document develops the Galerkin projection theory of the
Affine-DDPNM.  Starting from the broken-pressure
continuous-traction (C-DD) reference system, we prove:
(i)~the Schur solution at each enrichment level is the
$\Shatuu$-seminorm best approximation to the full C-DD trace solution;
(ii)~for nested trace subspaces the error admits an orthogonal
Pythagorean decomposition; and (iii)~the total volume error satisfies
a two-sided bound that separates the C-DD discretisation component
from the interface-modal-truncation component, with explicit
constants for the $A$-norm, $L^2$-norm, and broken-$H^1$ seminorm.

\textbf{Scope.}  The reference system throughout is the
broken-pressure C-DD formulation (discontinuous $P_1$ pressure across
pores).  Equivalence with the monolithic continuous-pressure
Taylor--Hood system is not established here; consequently the
discretisation term $\|\vel - \vel_h\|$ appearing in the separation
bounds is a C-DD discretisation error, not a standard FEM error.
All results assume planar interfaces with a fixed orthonormal frame.
\end{minipage}}
\end{center}

% ======================================================================
\section{C-DD reference system and trace subspaces}
% ======================================================================

\subsection{Notation and standing assumptions}

We work on a fixed finite-element mesh of the fluid domain
$\dom\subset\R^3$, partitioned into $N$ non-overlapping pore subdomains
$\{\pore{i}\}_{i=1}^N$ with $m$ planar internal interfaces
$\{\iface_k\}_{k=1}^m$.  On each pore we use the Taylor--Hood
$[P_2]^3$--$P_1$ pair with \textbf{discontinuous} pressure across
interfaces (the broken-pressure C-DD formulation).  All objects are
finite-dimensional; $h$ refers to the given mesh.

\begin{assumption}[Local algebraic stability]\label{ass:local}
  For every pore $i$:
  \begin{enumerate}[label=(\alph*),nosep,leftmargin=*]
    \item $A_i \in \R^{n_{u,i}\times n_{u,i}}$ is symmetric positive definite;
    \item $B_i \in \R^{n_{p,i}\times n_{u,i}}$ has full row rank;
    \item the constant function $1$ belongs to the local pressure space.
  \end{enumerate}
\end{assumption}

Assumption~\ref{ass:local}(a) is guaranteed when the pore has a
positive-measure no-slip wall segment.

\begin{assumption}[Matching interface trace spaces]\label{ass:matching}
  For every internal interface $\iface_k$ shared by pores $i$ and $j$,
  the two local velocity trace spaces coincide with a shared scalar
  nodal basis $\{\psi_{k,\ell}\}_{\ell=1}^{m_k}$ and a one-to-one DOF
  correspondence between the two sides.
\end{assumption}

\medskip\noindent
\textbf{Dimension conventions.}
$m_k$ = number of scalar trace DOFs on interface $k$;
$m_\Gamma = 3\sum_{k=1}^m m_k$ = total C-DD interface DOFs
(three components per trace node).
For pore $i$,
$r_i^u = 3\sum_{\gamma\in\mathcal{U}_i} m_{i,\gamma}$
(C-DD internal trace DOFs) and
$r_i^k = 3\sum_{\gamma\in\mathcal{K}_i} m_{i,\gamma}$
(C-DD boundary trace DOFs).
$n_i = 9m_i^u + m_i^k$ = affine primitive modes on pore $i$.

\subsection{Interface frames and affine scalar shapes}

On interface $\iface_k$ fix an orthonormal frame
$\{\normal_k,\bm{t}_{k,1},\bm{t}_{k,2}\}$ with $\normal_k$ directed
from the first pore of the pair to the second.  Let
$\bm{c}_k$ be the interface centroid and define the half-span-normalised
in-plane coordinates
\begin{equation}\label{eq:st}
  s_k(\bm{x}) = \frac{(\bm{x}-\bm{c}_k)\cdot\bm{t}_{k,1}}{a_k},\qquad
  t_k(\bm{x}) = \frac{(\bm{x}-\bm{c}_k)\cdot\bm{t}_{k,2}}{b_k},
\end{equation}
where $a_k,b_k > 0$ are the maximal absolute in-plane coordinate
values on $\iface_k$.  The three affine scalar shape functions are
\begin{equation}\label{eq:shapes}
  \varphi_0(\bm{x}) \equiv 1,\qquad
  \varphi_s(\bm{x}) = s_k(\bm{x}),\qquad
  \varphi_t(\bm{x}) = t_k(\bm{x}),
\end{equation}
and the mode set is
$\modeset := \{0,s,t\} \times \{n,t_1,t_2\}$.

\subsection{C-DD coupling, Schur complement, and energy seminorm}

For pore $i$ incident on $\iface_k$, define the oriented
inward-traction directions
$\bm{d}_{i,k,\beta}$ ($\beta\in\{n,t_1,t_2\}$) satisfying
$\bm{d}_{j,k,\beta} = -\bm{d}_{i,k,\beta}$ for the opposite pore $j$.

The C-DD coupling block for direction $\beta$ on $\iface_k$ is
\begin{equation}\label{eq:N}
  [N_{i,k}^{\beta}]_{\ell,j}
  = \int_{\iface_k} \psi_{k,\ell}\;
    \bm{d}_{i,k,\beta}\cdot\bm{\phi}_j^{(i)}\,\mathrm{d}S
  \;\in\; \R^{m_k\times n_{u,i}}.
\end{equation}
Collect the three-direction blocks on all internal ports of pore $i$
into $N_i^u \in \R^{r_i^u \times n_{u,i}}$, and on all boundary ports
into $N_i^k \in \R^{r_i^k \times n_{u,i}}$.  Let $\bm{\xi}^k \in
\R^{m_\Gamma^k}$ (with $m_\Gamma^k = \sum_i r_i^k$) denote the
globally prescribed boundary traction nodal values.  Through the
boundary affine restriction (see the companion
mathematical-properties document, Appendix~A) these are determined by
the prescribed inlet/outlet pressures.

The velocity response operator is
$H_i = A_i^{-1} - A_i^{-1}B_i^\top M_i^{-1}B_i A_i^{-1}$,
with $M_i = B_i A_i^{-1}B_i^\top$.
$H_i$ is symmetric positive semidefinite and satisfies
$H_i^\top A_i H_i = H_i$,
$\Ker(H_i) = \Range(B_i^\top)$.

The \textbf{internal-unknown block} of the local C-DD Schur complement is
\begin{equation}\label{eq:Shat-local}
  \Shatuu_i = N_i^u H_i (N_i^u)^\top \;\in\; \R^{r_i^u\times r_i^u}.
\end{equation}
Let $\widehat{R}_i^u : \R^{m_\Gamma} \to \R^{r_i^u}$ be the Boolean
restriction extracting the C-DD DOFs of pore $i$ from the global
vector of unknown interface DOFs.  The assembled C-DD Schur complement
(uu-block) is
\begin{equation}\label{eq:Shat}
  \Shatuu = \sum_{i=1}^N (\widehat{R}_i^u)^\top \Shatuu_i \widehat{R}_i^u
  \;\in\; \R^{m_\Gamma\times m_\Gamma}.
\end{equation}
$\Shatuu$ is symmetric positive semidefinite.

The C-DD right-hand side, with prescribed boundary tractions absorbed,
is
\begin{equation}\label{eq:Fhat}
  \Fhat = -\sum_{i=1}^N (\widehat{R}_i^u)^\top
    N_i^u H_i (N_i^k)^\top \widehat{R}_i^k \bm{\xi}^k
  \;\in\; \R^{m_\Gamma},
\end{equation}
where $\widehat{R}_i^k$ extracts the boundary DOFs of pore $i$ from
$\bm{\xi}^k$.

The \textbf{full C-DD system} for the unknown interface DOFs is
\begin{equation}\label{eq:full-sys}
  \Shatuu\,\bm{\xi} = \Fhat,\qquad \bm{\xi}\in\Vspace:=\R^{m_\Gamma}.
\end{equation}
Consistency of~\eqref{eq:full-sys} follows from well-posedness of the
underlying broken-pressure Stokes problem; we take it as given.
Define the energy seminorm on $\Vspace$ by
$\|\bm{\xi}\|_{\Shatuu}^2 := \bm{\xi}^\top \Shatuu \bm{\xi}$.

\subsection{Restriction operators and nested trace subspaces}

\begin{definition}[Per-interface restriction matrices]\label{def:Pkl}
  For each interface $\iface_k$, define the $L^2$-projection vectors
  $p_k^\alpha = M_k^{-1} b_k^\alpha \in \R^{m_k}$ for
  $\alpha\in\{0,s,t\}$, where
  $M_k$ is the face mass matrix on $\iface_k$ and
  $[b_k^\alpha]_\ell = \int_{\iface_k} \varphi_\alpha\psi_{k,\ell}\,\mathrm{d}S$.
  For level $\ell\in\{0,1,2\}$ with $d_\ell$ coefficient DOFs per interface,
  \begin{itemize}
    \item $\ell=0$ (classic P0, $d_0=1$):
      $P_{k,0} = \begin{bmatrix} p_k^0 \\ \zero \\ \zero \end{bmatrix}
       \in \R^{3m_k\times 1}$;
    \item $\ell=1$ (constant-vector P0, $d_1=3$):
      $P_{k,1} = \begin{bmatrix}
        p_k^0 & \zero   & \zero   \\
        \zero & p_k^0   & \zero   \\
        \zero & \zero   & p_k^0
      \end{bmatrix} \in \R^{3m_k\times 3}$;
    \item $\ell=2$ (full affine, $d_2=9$):
      $P_{k,2} = \begin{bmatrix}
        p_k^0 & \zero & \zero & p_k^s & \zero & \zero & p_k^t & \zero & \zero \\[2pt]
        \zero & p_k^0 & \zero & \zero & p_k^s & \zero & \zero & p_k^t & \zero \\[2pt]
        \zero & \zero & p_k^0 & \zero & \zero & p_k^s & \zero & \zero & p_k^t
      \end{bmatrix} \in \R^{3m_k\times 9}$.
  \end{itemize}
  Each column occupies only the block-row matching its direction
  $\beta\in\{n,t_1,t_2\}$.
\end{definition}

The global restriction at level $\ell$ is the block-diagonal assembly
\begin{equation}\label{eq:Pcal}
  \Pcal_\ell := \bigoplus_{k=1}^m P_{k,\ell}
  \;:\; \R^{d_\ell m} \;\longrightarrow\; \Vspace.
\end{equation}
The corresponding \textbf{trace subspace} is
$\Wspace_\ell := \Range(\Pcal_\ell) \subset \Vspace$.

\begin{lemma}[Nested trace subspaces]\label{lem:nested}
  $\Wspace_0 \subset \Wspace_1 \subset \Wspace_2 \subset \Vspace$.
\end{lemma}

\begin{proof}
  Each $P_{k,0}$ (resp.\ $P_{k,1}$) is obtained from $P_{k,1}$
  (resp.\ $P_{k,2}$) by selecting a subset of its columns.
  The block-diagonal structure of $\Pcal_\ell$ preserves the inclusion
  of ranges.
\end{proof}

\begin{remark}[Coefficient-space nesting is not literal]\label{rem:coeff-nest}
  The coefficient spaces $\R^{d_0 m}$, $\R^{d_1 m}$, $\R^{d_2 m}$ are
  Euclidean spaces of different dimensions and cannot be written as
  subsets of one another.  The genuine nested structure lives at the
  level of \emph{trace subspaces} $\Wspace_\ell \subset \Vspace$.
  All results below are stated in terms of these trace subspaces.
  When we refer to ``level $\ell$'' we mean the subspace $\Wspace_\ell$
  and its associated restriction operator $\Pcal_\ell$.
\end{remark}

% ======================================================================
\section{Best approximation in the energy seminorm}
% ======================================================================

For a fixed level $\ell$, the \textbf{reduced Schur complement} and
\textbf{reduced right-hand side} are
\begin{equation}\label{eq:S-ell}
  \Schur_\ell = \Pcal_\ell^\top \Shatuu\, \Pcal_\ell
  \;\in\; \R^{d_\ell m \times d_\ell m},\qquad
  \bm{F}_\ell = \Pcal_\ell^\top \Fhat.
\end{equation}

\begin{lemma}[Reduced consistency]\label{lem:reduced-cons}
  If the full C-DD system~\eqref{eq:full-sys} is consistent, then for
  every level $\ell$,
  $\bm{F}_\ell \in \Range(\Schur_\ell)$.
\end{lemma}

\begin{proof}
  Let $\bm{\xi}$ be any solution of $\Shatuu\,\bm{\xi} = \Fhat$.
  Set $A_\ell = (\Shatuu)^{1/2}\Pcal_\ell$; then
  $\Schur_\ell = A_\ell^\top A_\ell$ and
  $\bm{F}_\ell = A_\ell^\top (\Shatuu)^{1/2}\bm{\xi}$.
  Finite-dimensional linear algebra gives
  $\Range(A_\ell^\top A_\ell) = \Range(A_\ell^\top)$,
  hence $\bm{F}_\ell \in \Range(\Schur_\ell)$.
\end{proof}

Consequently there exists at least one solution $\trac_\ell \in
\R^{d_\ell m}$ of the reduced system
$\Schur_\ell \trac_\ell = \bm{F}_\ell$.
Define the corresponding C-DD trace approximation
\begin{equation}\label{eq:xi-ell}
  \bm{\xi}_\ell := \Pcal_\ell \trac_\ell \;\in\; \Wspace_\ell \subset \Vspace.
\end{equation}
Let $\bm{\xi} \in \Vspace$ be any fixed solution of the full C-DD
system $\Shatuu\,\bm{\xi} = \Fhat$.

\begin{theorem}[Galerkin best approximation]\label{thm:best}
  Under Assumptions~\ref{ass:local} and~\ref{ass:matching},
  for each level $\ell\in\{0,1,2\}$,
  \begin{equation}\label{eq:best}
    \|\bm{\xi} - \bm{\xi}_\ell\|_{\Shatuu}
    = \min_{\bm{\mu}\in\R^{d_\ell m}}
      \|\bm{\xi} - \Pcal_\ell\bm{\mu}\|_{\Shatuu}.
  \end{equation}
  When $\Shatuu$ is semidefinite the minimiser need not be unique,
  but the seminorm value is independent of the choice of
  representatives for both $\bm{\xi}$ and $\trac_\ell$.
\end{theorem}

\begin{proof}
  From $\Shatuu\,\bm{\xi} = \Fhat$,
  $\Schur_\ell \trac_\ell = \bm{F}_\ell$,
  $\Schur_\ell = \Pcal_\ell^\top \Shatuu\,\Pcal_\ell$, and
  $\bm{F}_\ell = \Pcal_\ell^\top \Fhat$, we obtain
  \begin{equation}\label{eq:ortho}
    \Pcal_\ell^\top \Shatuu\,(\bm{\xi}_\ell - \bm{\xi}) = \zero.
  \end{equation}
  Hence $\bm{e}_\ell := \bm{\xi} - \bm{\xi}_\ell$ is
  $\Shatuu$-orthogonal to $\Wspace_\ell = \Range(\Pcal_\ell)$.

  For arbitrary $\bm{\mu}\in\R^{d_\ell m}$, decompose
  $\bm{\xi} - \Pcal_\ell\bm{\mu}
   = \bm{e}_\ell + \bm{v}$
  with $\bm{v} := \bm{\xi}_\ell - \Pcal_\ell\bm{\mu}
  \in \Wspace_\ell$.
  By~\eqref{eq:ortho}, $\bm{e}_\ell^\top \Shatuu\,\bm{v} = 0$, and the
  Pythagorean identity in the $\Shatuu$ seminorm yields
  $\|\bm{\xi} - \Pcal_\ell\bm{\mu}\|_{\Shatuu}^2
   = \|\bm{e}_\ell\|_{\Shatuu}^2 + \|\bm{v}\|_{\Shatuu}^2
   \ge \|\bm{e}_\ell\|_{\Shatuu}^2$.
  Equality at $\bm{\mu} = \trac_\ell$ proves~\eqref{eq:best}.

  Semidefinite independence: if $\trac_\ell,\trac_\ell'$ both solve
  $\Schur_\ell\trac = \bm{F}_\ell$, then
  $\Pcal_\ell(\trac_\ell-\trac_\ell') \in \Ker(\Shatuu)$,
  so the seminorm on both sides of~\eqref{eq:best} is unchanged.
  Similarly, any two full solutions differ by an element of
  $\Ker(\Shatuu)$, which leaves the seminorm unchanged
  (since $\bm{\xi}_\ell \in \Wspace_\ell$ and
   $\Pcal_\ell^\top\Shatuu(\bm{\xi}_\ell-\bm{\xi})=\zero$).
\end{proof}

% ======================================================================
\section{Nested orthogonal decomposition}
% ======================================================================

\begin{theorem}[Orthogonal decomposition for nested trace subspaces]\label{thm:nested}
  Under the same assumptions as Theorem~\ref{thm:best}, let
  $\ell < \ell'$, so that $\Wspace_\ell \subset \Wspace_{\ell'}$ by
  Lemma~\ref{lem:nested}.  Then
  \begin{equation}\label{eq:nested}
    \|\bm{\xi} - \bm{\xi}_\ell\|_{\Shatuu}^2
    = \|\bm{\xi} - \bm{\xi}_{\ell'}\|_{\Shatuu}^2
      + \|\bm{\xi}_{\ell'} - \bm{\xi}_\ell\|_{\Shatuu}^2.
  \end{equation}
\end{theorem}

\begin{proof}
  By Theorem~\ref{thm:best} at level $\ell'$,
  $\bm{e}_{\ell'} = \bm{\xi} - \bm{\xi}_{\ell'}$ is
  $\Shatuu$-orthogonal to $\Wspace_{\ell'}$.
  The difference $\bm{\xi}_{\ell'} - \bm{\xi}_\ell$ satisfies
  $\bm{\xi}_{\ell'} \in \Wspace_{\ell'}$ and
  $\bm{\xi}_\ell \in \Wspace_\ell \subset \Wspace_{\ell'}$,
  hence $\bm{\xi}_{\ell'} - \bm{\xi}_\ell \in \Wspace_{\ell'}$.
  Therefore $\bm{e}_{\ell'} \perp_{\Shatuu}
  (\bm{\xi}_{\ell'} - \bm{\xi}_\ell)$.
  The decomposition
  $\bm{\xi} - \bm{\xi}_\ell = \bm{e}_{\ell'} +
   (\bm{\xi}_{\ell'} - \bm{\xi}_\ell)$
  is $\Shatuu$-orthogonal; the Pythagorean identity gives~\eqref{eq:nested}.
\end{proof}

\begin{corollary}[Three-level gain decomposition]\label{cor:gains}
  For $\Wspace_0 \subset \Wspace_1 \subset \Wspace_2$,
  \begin{equation}\label{eq:gains}
    \|\bm{\xi} - \bm{\xi}_0\|_{\Shatuu}^2
    = \|\bm{\xi} - \bm{\xi}_2\|_{\Shatuu}^2
      + \|\bm{\xi}_2 - \bm{\xi}_1\|_{\Shatuu}^2
      + \|\bm{\xi}_1 - \bm{\xi}_0\|_{\Shatuu}^2.
  \end{equation}
  Interpretation of the three right-hand-side terms:
  \begin{itemize}
    \item $\|\bm{\xi} - \bm{\xi}_2\|_{\Shatuu}^2$:
      \textbf{affine residual} --- error after the full nine-mode
      affine approximation;
    \item $\|\bm{\xi}_2 - \bm{\xi}_1\|_{\Shatuu}^2$:
      \textbf{linear gain} --- reduction from adding the six linear
      modes $\{s,t\}\times\{n,t_1,t_2\}$ to the constant-vector space;
    \item $\|\bm{\xi}_1 - \bm{\xi}_0\|_{\Shatuu}^2$:
      \textbf{tangential gain} --- reduction from adding the two
      tangential constant modes to the single normal constant mode.
  \end{itemize}
\end{corollary}

\begin{proof}
  Apply Theorem~\ref{thm:nested} with $(\ell,\ell')=(0,1)$ and
  $(\ell,\ell')=(1,2)$, then substitute.
\end{proof}

\begin{remark}[Four-space causal decomposition]\label{rem:four-space}
  The three-level chain $\Wspace_0\subset\Wspace_1\subset\Wspace_2$
  lumps together the six linear modes as a single gain term and cannot
  distinguish whether the dominant missing physics is normal-direction
  linear variation or tangential-direction linear variation.
  To resolve this one needs four trace subspaces defined by their
  per-interface mode content:
  \begin{align*}
    \Wspace_{0n} &= \operatorname{span}\{1\cdot n\}
      \quad\text{(1 mode/interface: classic P0)},\\
    \Wspace_{0v} &= \operatorname{span}\{1\cdot n,\;1\cdot t_1,\;1\cdot t_2\}
      \quad\text{(3 modes/interface: constant vector)},\\
    \Wspace_{1n} &= \operatorname{span}\{1\cdot n,\;s\cdot n,\;t\cdot n\}
      \quad\text{(3 modes/interface: linear normal only)},\\
    \Wspace_{1v} &= \operatorname{span}\bigl(\{1,s,t\}\times\{n,t_1,t_2\}\bigr)
      \quad\text{(9 modes/interface: full affine)}.
  \end{align*}
  These four spaces form a commutative diamond:
  $\Wspace_{0n}\subset\Wspace_{0v}\subset\Wspace_{1v}$ \emph{and}
  $\Wspace_{0n}\subset\Wspace_{1n}\subset\Wspace_{1v}$,
  but $\Wspace_{0v}\not\subset\Wspace_{1n}$ and
  $\Wspace_{1n}\not\subset\Wspace_{0v}$.
  Pythagorean decompositions along the two distinct nested paths give
  two orthogonal partitions of the classic-P0 error:
  \[
  \begin{aligned}
  \|\bm{\xi}-\bm{\xi}_{0n}\|_{\Shatuu}^2
  &= \|\bm{\xi}-\bm{\xi}_{1v}\|_{\Shatuu}^2
   + \|\bm{\xi}_{1v}-\bm{\xi}_{0v}\|_{\Shatuu}^2
   + \|\bm{\xi}_{0v}-\bm{\xi}_{0n}\|_{\Shatuu}^2,\\
  &= \|\bm{\xi}-\bm{\xi}_{1v}\|_{\Shatuu}^2
   + \|\bm{\xi}_{1v}-\bm{\xi}_{1n}\|_{\Shatuu}^2
   + \|\bm{\xi}_{1n}-\bm{\xi}_{0n}\|_{\Shatuu}^2.
  \end{aligned}
  \]
  The first path asks: after adding all three constant-vector modes
  (tangential gain), how much more does full affine add (linear gain)?
  The second path asks: after adding linear normal modes
  (normal-direction enrichment), how much more do tangential modes add?
  Comparing the per-interface contributions from the two paths resolves
  whether the error is dominated by missing normal-direction variation
  or missing tangential traction.  This extension to four spaces and the
  associated two-path decomposition is left to future work.
\end{remark}

% ======================================================================
\section{Velocity reconstruction and trace error identity}
% ======================================================================

\subsection{Velocity reconstruction from C-DD trace data}

For a C-DD trace vector $\bm{\zeta}\in\Vspace$ (unknown internal DOFs)
and prescribed boundary traction data $\bm{\xi}^k$, the local velocity
field on pore $i$ is
\begin{equation}\label{eq:reconstruct}
  U_i(\bm{\zeta};\,\bm{\xi}^k)
  = -H_i\Bigl[(N_i^u)^\top \widehat{R}_i^u \bm{\zeta}
               + (N_i^k)^\top \widehat{R}_i^k \bm{\xi}^k\Bigr]
  \;\in\; \R^{n_{u,i}}.
\end{equation}
The first term couples the unknown interface DOFs to the pore
interior; the second term injects the prescribed boundary tractions.
The global reconstructed velocity field $\vel_h(\bm{\zeta})$ is the
discontinuous assembly of the local fields $\{U_i(\bm{\zeta};\,
\bm{\xi}^k)\}$ across pores.

\begin{lemma}[Trace-to-velocity error identity]\label{lem:vel-id}
  For any two trace vectors $\bm{\zeta},\bm{\zeta}'\in\Vspace$ sharing
  the same prescribed boundary data $\bm{\xi}^k$,
  \begin{equation}\label{eq:vel-id}
    \sum_{i=1}^N \|U_i(\bm{\zeta}) - U_i(\bm{\zeta}')\|_{A_i}^2
    = \|\bm{\zeta} - \bm{\zeta}'\|_{\Shatuu}^2,
  \end{equation}
  where $\|U\|_{A_i}^2 := U^\top A_i U$.
\end{lemma}

\begin{proof}
  The boundary coupling terms $(N_i^k)^\top \widehat{R}_i^k \bm{\xi}^k$
  are identical for both arguments and cancel in the difference:
  \begin{align*}
    U_i(\bm{\zeta}) - U_i(\bm{\zeta}')
    &= -H_i (N_i^u)^\top \widehat{R}_i^u (\bm{\zeta} - \bm{\zeta}').
  \end{align*}
  Using $H_i^\top A_i H_i = H_i$,
  \begin{align*}
    \|U_i(\bm{\zeta}) - U_i(\bm{\zeta}')\|_{A_i}^2
    &= (\bm{\zeta}-\bm{\zeta}')^\top \widehat{R}_i^{u\top}
       N_i^u H_i A_i H_i (N_i^u)^\top \widehat{R}_i^u
       (\bm{\zeta}-\bm{\zeta}') \\
    &= (\bm{\zeta}-\bm{\zeta}')^\top \widehat{R}_i^{u\top}
       \Shatuu_i \widehat{R}_i^u (\bm{\zeta}-\bm{\zeta}').
  \end{align*}
  Summing over $i$ and using~\eqref{eq:Shat} gives~\eqref{eq:vel-id}.
\end{proof}

\subsection{Notation for the true solution and its approximations}

Let $\vel$ denote the exact PDE solution of the Stokes problem.
Let $\vel_h := \vel_h(\bm{\xi})$ be the C-DD full-trace reconstruction
using a solution $\bm{\xi}$ of~\eqref{eq:full-sys}.
For level $\ell$, let $\vel_\ell := \vel_h(\bm{\xi}_\ell)$.
All three velocities share the same prescribed boundary tractions
$\bm{\xi}^k$.

Define the following broken norms on the product velocity space
$\prod_i H^1(\pore{i})^3$.  For a (possibly discontinuous) velocity
field $\bm{v}$,
\begin{align}
  \|\bm{v}\|_A^2 &:= \sum_{i=1}^N \nu\int_{\pore{i}}
                     \nabla\bm{v}:\nabla\bm{v}\,\mathrm{d}x,
                     \label{eq:A-norm} \\[4pt]
  \|\bm{v}\|_{L^2}^2 &:= \sum_{i=1}^N \int_{\pore{i}}
                         |\bm{v}|^2\,\mathrm{d}x, \label{eq:L2-norm} \\[4pt]
  |\bm{v}|_{H^1}^2 &:= \sum_{i=1}^N \int_{\pore{i}}
                        \nabla\bm{v}:\nabla\bm{v}\,\mathrm{d}x
                        = \frac{1}{\nu}\|\bm{v}\|_A^2. \label{eq:H1-semi}
\end{align}
For a discrete function $\bm{v}_h$ with local coefficient vectors
$\{V_i\}$, the $A$-norm satisfies
$\|\bm{v}_h\|_A^2 = \sum_i V_i^\top A_i V_i$.

% ======================================================================
\section{Volume error---trace error separation}
% ======================================================================

\subsection{Two-sided bound: the true error floor}

Set
\begin{equation}\label{eq:DE}
  D_h := \|\vel - \vel_h\|_A,\qquad
  E_\ell := \|\vel_h - \vel_\ell\|_A
          = \|\bm{\xi} - \bm{\xi}_\ell\|_{\Shatuu},
\end{equation}
where the second equality is Lemma~\ref{lem:vel-id} with
$\bm{\zeta}=\bm{\xi}$, $\bm{\zeta}'=\bm{\xi}_\ell$.

\begin{theorem}[Two-sided volume-trace bound: $A$-norm]\label{thm:sep-A}
  Under Assumptions~\ref{ass:local} and~\ref{ass:matching},
  for each level $\ell$,
  \begin{equation}\label{eq:sep-A}
    |E_\ell - D_h|
    \;\le\; \|\vel - \vel_\ell\|_A
    \;\le\; E_\ell + D_h.
  \end{equation}
  In particular, when the C-DD discretisation error is negligible
  ($D_h \ll E_\ell$), the total error is essentially the trace error:
  $\|\vel - \vel_\ell\|_A = E_\ell + O(D_h)$.
\end{theorem}

\begin{proof}
  By the reverse and forward triangle inequalities for the $A$-norm,
  \[
    \bigl|\|\vel_h - \vel_\ell\|_A - \|\vel - \vel_h\|_A\bigr|
    \le \|\vel - \vel_\ell\|_A
    \le \|\vel_h - \vel_\ell\|_A + \|\vel - \vel_h\|_A.
  \]
  Substituting $E_\ell$ and $D_h$ from~\eqref{eq:DE} gives the result.
\end{proof}

\begin{remark}[Why this is an error floor --- and what it does not prove]\label{rem:floor}
  On a fixed mesh $h$, the lower bound $E_\ell - D_h$ shows that the
  total error cannot drop below the trace-modal-truncation error
  $E_\ell$ by more than the C-DD discretisation error $D_h$.
  When $D_h \ll E_\ell$, the error on this mesh is trace-dominated:
  $\|\vel-\vel_\ell\|_A = E_\ell + O(D_h)$, and further mesh refinement
  on this particular grid would not help---the limiting factor is the
  interface subspace $\Wspace_\ell$.

  \textbf{Caution.}  This is a \emph{single-mesh} statement.
  Both $E_\ell = E_{\ell,h}$ and $D_h$ depend on $h$.  A mesh-independent
  asymptotic error floor additionally requires
  $\liminf_{h\to 0} E_{\ell,h} > 0$ (or convergence to a positive
  limit $E_{\ell}^* > 0$), which is not proved here.  If $E_{\ell,h}$
  were to decay with mesh refinement (e.g.\ because the trace space
  becomes richer on finer meshes), the total error could still decrease.
  The numerical observation that errors plateau at $\sim$6.5\% (planar)
  and $\sim$14\% (watershed) under mesh refinement is empirical evidence
  for a positive asymptotic floor; its proof requires analysing the
  $h$-dependence of the C-DD trace space and the restriction operator
  $\Pcal_\ell$.

  \textbf{Computability.}  $E_\ell = \|\bm{\xi}-\bm{\xi}_\ell\|_{\Shatuu}$
  is computable from the full and reduced discrete solutions.
  $D_h = \|\vel-\vel_h\|_A$ requires the exact PDE solution $\vel$ and is
  \emph{not} directly computable; it must be estimated via a manufactured
  solution, an overkill reference computation, or a separate discretisation
  error estimator.
\end{remark}

\subsection{Analogous bounds in $L^2$ and broken $H^1$}

\begin{theorem}[Two-sided bound: $L^2$-norm]\label{thm:sep-L2}
  Under the same assumptions, suppose each pore $\pore{i}$ admits a
  Poincar\'e--Friedrichs constant $C_{P,i}>0$ such that
  $\|\bm{v}\|_{L^2(\pore{i})} \le C_{P,i}|\bm{v}|_{H^1(\pore{i})}$
  for all $\bm{v}$ vanishing on the wall portion of
  $\partial\pore{i}$.  Set $C_P := \max_i C_{P,i}$ and define
  \begin{equation}\label{eq:TL2}
    T_\ell^{L^2} := \|\vel_h - \vel_\ell\|_{L^2},\qquad
    D_h^{L^2} := \|\vel - \vel_h\|_{L^2}.
  \end{equation}
  Then the following hold.
  \begin{enumerate}[label=(\alph*),nosep,leftmargin=*]
    \item \textbf{Two-sided bound with $T_\ell^{L^2}$:}
      $|T_\ell^{L^2} - D_h^{L^2}|
       \le \|\vel - \vel_\ell\|_{L^2}
       \le T_\ell^{L^2} + D_h^{L^2}$.
    \item \textbf{Poincar\'e upper bound:}
      $T_\ell^{L^2} \le \frac{C_P}{\sqrt{\nu}}\,E_\ell$, and therefore
      $\|\vel - \vel_\ell\|_{L^2}
       \le D_h^{L^2} + \frac{C_P}{\sqrt{\nu}}\,E_\ell$.
  \end{enumerate}
  No corresponding lower bound in terms of $E_\ell$ is available without
  a reverse stability estimate (e.g.\ an inverse inequality bounding
  the $A$-norm from above by the $L^2$-norm on the discrete velocity
  space).
\end{theorem}

\begin{proof}
  (a)~Apply the reverse and forward triangle inequalities for the
  $L^2$-norm to $\vel-\vel_\ell = (\vel-\vel_h) + (\vel_h-\vel_\ell)$.
  (b)~From Lemma~\ref{lem:vel-id} and per-pore Poincar\'e--Friedrichs,
  $(T_\ell^{L^2})^2
   = \sum_i \|U_i(\bm{\xi})-U_i(\bm{\xi}_\ell)\|_{L^2(\pore{i})}^2
   \le \frac{C_P^2}{\nu}\sum_i \|U_i(\bm{\xi})-U_i(\bm{\xi}_\ell)\|_{A_i}^2
   = \frac{C_P^2}{\nu}\,E_\ell^2$.
  Taking square roots and substituting into the upper bound of (a)
  gives the second statement.
\end{proof}

\begin{theorem}[Two-sided bound: broken-$H^1$ seminorm]\label{thm:sep-H1}
  Under the same assumptions, define
  $T_\ell^{H^1} := |\vel_h - \vel_\ell|_{H^1}$ and
  $D_h^{H^1} := |\vel - \vel_h|_{H^1}$.
  From~\eqref{eq:H1-semi} and Lemma~\ref{lem:vel-id},
  $T_\ell^{H^1} = E_\ell/\sqrt{\nu}$ exactly.
  Then
  \begin{equation}\label{eq:sep-H1}
    |T_\ell^{H^1} - D_h^{H^1}|
    \;\le\; |\vel - \vel_\ell|_{H^1}
    \;\le\; T_\ell^{H^1} + D_h^{H^1}
    \;=\; \frac{1}{\sqrt{\nu}}\,E_\ell + D_h^{H^1}.
  \end{equation}
\end{theorem}

\begin{proof}
  The identity $T_\ell^{H^1} = E_\ell/\sqrt{\nu}$ follows from
  $|\bm{v}|_{H^1}^2 = \|\bm{v}\|_A^2/\nu$ and Lemma~\ref{lem:vel-id}.
  The seminorm triangle inequalities give the two-sided bound.
  Unlike the $L^2$ case, both the upper and the lower bound are
  available here because $T_\ell^{H^1}$ is exactly computable
  from $E_\ell$ (no additional inequality needed).
\end{proof}

\begin{remark}[Pessimism of the $L^2$ constant]\label{rem:CP}
  The constant $C_P = \max_i C_{P,i}$ is the worst-case Poincar\'e
  constant across all pores and can be very pessimistic for irregular
  watershed geometries.  In practice the effective constant for the
  specific error field $\vel_h-\vel_\ell$ is given by the Rayleigh
  quotient $\|\vel_h-\vel_\ell\|_{L^2} / \|\vel_h-\vel_\ell\|_A$
  and can be evaluated numerically.
\end{remark}

% ======================================================================
\section{Discussion}
% ======================================================================

\subsection{Summary of the theoretical chain}

\begin{enumerate}[leftmargin=*,itemsep=4pt]
  \item \textbf{Best approximation} (Theorem~\ref{thm:best}):
    the level-$\ell$ DDPNM trace solution is the $\Shatuu$-seminorm
    projection of the full C-DD trace solution onto $\Wspace_\ell$.
    This is a standard Galerkin optimality result.
  \item \textbf{Nested decomposition} (Theorem~\ref{thm:nested},
    Corollary~\ref{cor:gains}):
    error reductions obtained by enriching the trace subspace are
    $\Shatuu$-orthogonal.  The total classic-P0 error decomposes into
    the affine residual plus two orthogonal gain terms.
  \item \textbf{Error floor} (Theorems~\ref{thm:sep-A}--\ref{thm:sep-H1}):
    the total volume error satisfies a two-sided bound
    $|E_\ell - D_h| \le \|\vel-\vel_\ell\| \le E_\ell + D_h$
    in each of the three norms ($A$, $L^2$, broken-$H^1$).
    When $D_h \ll E_\ell$, the trace error $E_\ell$ is the effective
    error floor.
\end{enumerate}

\subsection{Limitations of the current theory and directions for adaptive enrichment}

The orthogonal gain decomposition (Corollary~\ref{cor:gains})
provides the algebraic foundation for residual-driven adaptive
enrichment.  The enrichment algorithm in
\texttt{ddpnm3d/hierarchy.py}
(\texttt{run\_adaptive\_hierarchy}) implements a two-phase loop:
starting from Level~0 on all interfaces, residual indicators
(from \texttt{ddpnm\_core/estimate.py}) drive D\"orfler
  marking to
upgrade selected interfaces to Level~1, then Level~2.

\textbf{What the algorithm converges to.}
The enrichment operates on a fixed mesh $h$ with a fixed set of
three trace subspaces $\Wspace_0\subset\Wspace_1\subset\Wspace_2$.
Because the state space is finite ($3^m$ possible level assignments),
the algorithm terminates in at most $2m$ steps; the limit is the
all-Level-2 solution $\bm{\xi}_2$, with error floor
$E_2 = \|\bm{\xi}-\bm{\xi}_2\|_{\Shatuu}$.
The total error \emph{cannot} reach zero for three independent reasons:
\begin{enumerate}[nosep,leftmargin=*]
  \item \textbf{Modal truncation:} the nine affine modes per interface
    span a strict subspace of the full C-DD trace space (${\sim}100$
    nodal DOFs per interface in 3D); the complement is never enriched.
  \item \textbf{C-DD $\neq$ PDE:} the reference system uses broken
    pressure; even the full C-DD solution $\vel_h$ does not equal the
    monolithic Taylor--Hood solution (let alone the PDE solution $\vel$).
  \item \textbf{Geometric error (watershed only):} per-facet normal
    dispersion on faceted interfaces introduces an additional
    representation gap that the fixed-frame affine space cannot close.
\end{enumerate}
Overcoming these would require, respectively, higher-order modal
enrichment ($P_2$, $P_3$, \dots), a proof of C-DD--monolithic
equivalence (or a monolithic reference theory), and per-facet normal
frames for non-planar interfaces.

\textbf{Toward a contraction theorem.}
For a D\"orfler-driven loop to provably reduce the error by a fixed
factor per iteration ($E_{k+1}\le q E_k$, $q<1$), the following
ingredients are needed beyond what is established here:
\begin{enumerate}[nosep,leftmargin=*]
  \item \textbf{Reliability:} each residual indicator $\eta_i$ must
    satisfy $\eta_i^2 \ge c_{\rm rel}\,g_i^2$, where $g_i^2 =
    \|\bm{\xi}_{\rm enriched}-\bm{\xi}_{\rm current}\|_{\Shatuu}^2$
    restricted to interface $i$ (the local contribution to the
    orthogonal gain term in Corollary~\ref{cor:gains}).
  \item \textbf{Efficiency:} $\eta_i^2 \le C_{\rm eff}\,g_i^2$ (the
    indicator does not over-estimate the available gain).
  \item \textbf{Estimator reduction} or a \textbf{saturation} property
    linking the marked set to the global error reduction.
\end{enumerate}
None of these are proved in the present work; the indicators in
\texttt{estimate.py} are currently validated only empirically.

\textbf{Summary of what is and is not proved.}
Theorem~\ref{thm:best} (best approximation) and
Theorem~\ref{thm:nested} (orthogonal decomposition) are rigorous
within the C-DD framework.  Theorem~\ref{thm:sep-A}--\ref{thm:sep-H1}
provide two-sided volume-trace separation bounds that define the error
floor on a fixed mesh.  What is missing for a full adaptive convergence
theory is: reliability/efficiency of the indicators, a contraction or
estimator-reduction theorem, the four-space diamond decomposition
(Remark~\ref{rem:four-space}), and the bridge from the C-DD reference
system to either the monolithic FEM or the PDE.

\subsection{What is not yet proved}

\begin{enumerate}[leftmargin=*,itemsep=4pt]
  \item \textbf{C-DD $\neq$ monolithic FEM.}
    The C-DD reference system uses discontinuous pressure across
    pores; the monolithic Taylor--Hood system uses globally
    continuous $P_1$ pressure.  Equivalence of the two velocity
    solutions is not established.  Consequently the discretisation
    error $D_h$ appearing in the separation bounds is a C-DD
    discretisation error, not directly comparable to standard FE
    error estimates.
  \item \textbf{Mesh-independent error floor.}
    The two-sided bounds in Theorems~\ref{thm:sep-A}--\ref{thm:sep-H1}
    hold on a fixed mesh.  Proving that $E_{\ell,h}$ remains bounded
    away from zero as $h\to 0$ (the asymptotic floor observed
    numerically at $\sim$6.5\% and $\sim$14\%) requires an analysis of
    the $h$-dependence of the trace space and the restriction operator
    $\Pcal_\ell$ (Remark~\ref{rem:floor}).
  \item \textbf{Reliability and efficiency of residual indicators.}
    The link between the orthogonal gain terms and the per-interface
    residual indicators computed in \texttt{ddpnm\_core/estimate.py}
    is currently empirical.  Upper and lower bounds (efficiency and
    reliability \`a la Babu\v{s}ka--Rheinboldt) are required for a
    contraction theorem.
  \item \textbf{Contraction theorem for adaptive enrichment.}
    Even with reliable and efficient indicators, a D\"orfler marking
    strategy additionally requires an estimator-reduction or saturation
    property to guarantee $E_{k+1}\le q E_k$ with $q<1$.
  \item \textbf{Four-space causal decomposition.}
    The three-level chain $\Wspace_0\subset\Wspace_1\subset\Wspace_2$
    does not separate normal-direction linear enrichment from
    tangential-direction linear enrichment.  The four-space diamond
    (Remark~\ref{rem:four-space}) with its two-path Pythagorean
    decomposition is defined but not yet implemented or numerically
    evaluated.
  \item \textbf{Non-planar interfaces.}
    All results assume planar interfaces with a fixed orthonormal
    frame.  Watershed partitions produce faceted interfaces whose
    per-facet normal variation introduces additional geometric error
    not captured by the present analysis.
  \item \textbf{Higher-order modal enrichment.}
    Overcoming the modal truncation error floor would require enriching
    beyond the nine affine modes (e.g.\ $P_2$ polynomial shapes
    $\{1,s,t,s^2,st,t^2\}\times\{n,t_1,t_2\}=18$ modes/interface).
    The nested decomposition theory extends naturally to any finite
    chain of nested trace subspaces, but the construction and numerical
    validation of $P_2$ restriction operators is left to future work.
\end{enumerate}

% ======================================================================
\begin{thebibliography}{9}

\bibitem{sun2025ddpnm}
S.~Sun, Z.~Wang, L.~Zhang, J.~Zhao.
\newblock A domain decomposition approach to pore-network modeling of porous
  media flow.
\newblock \textit{arXiv:2510.13429v2}, 2026.

\bibitem{toselli2005}
A.~Toselli and O.~Widlund.
\newblock \textit{Domain Decomposition Methods---Algorithms and Theory}.
\newblock Springer, 2005.

\bibitem{horn}
R.~A.~Horn and C.~R.~Johnson.
\newblock \textit{Matrix Analysis}, 2nd ed.
\newblock Cambridge University Press, 2012.

\end{thebibliography}

\end{document}
